\section{Objects tree}\label{objects-tree}\index{objects tree}

The \emph{Objects tree} is the modern way of building a scene in Mandelbulber. Instead of editing
a single flat list of fractal parameters, every element of the image (a fractal formula, a geometric
primitive, or a group that combines several objects) is a \emph{node} in a tree, and the scene is
assembled from those nodes like in a typical 3D application. This chapter describes the
\emph{Objects tree} interface: what each control does and how to use it in practice.

The Objects tree lives in the \emph{Objects} dock -- the settings pane on the right side of the
main window -- on the tab titled \emph{Objects tree}. If the dock is not visible, the standard
layout can be restored with \emph{View / Reset window layout to default}. Older scene files
(\emph{.fract}) are converted automatically into an object tree when they are loaded, so previously
created scenes keep working without any action from the user.

\subsection{Window layout}\label{objects-tree-layout}\index{objects tree!layout}

The tab consists of three parts (figure \ref{objects_tree_overview}):

\begin{itemize}
	\item a row of \textbf{buttons} at the top: \emph{Add group}, \emph{Add Fractal},
	\emph{Add Primitive}, \emph{Duplicate}, \emph{Delete object}, \emph{Refresh}, and the
	\emph{Display wire-frame previews} checkbox;
	\item the \textbf{tree view} listing all objects of the scene;
	\item the \textbf{editor panel} below the tree, showing the parameters of the currently
	selected object. The splitter between the tree and the editor can be dragged to give more
	space to either part.
\end{itemize}

Above the tab there is the standard row of local buttons (\emph{Load}, \emph{Save}, \emph{Reset},
\emph{Randomize}, \emph{Navi}) which works the same way as on the other settings panes.

\simpleImageWithCaptionThirdWidth{img/manual/media/objects_tree_dock_overview.png}
{The Objects tree dock: buttons, tree and editor panel}
{objects_tree_overview}{H}

\subsection{The tree and its columns}\label{objects-tree-columns}\index{objects tree!columns}

Every node is displayed as one line in the tree. The tree always has a single root called
\emph{World}, which cannot be renamed or deleted; all scene objects hang from it, directly or
inside groups. The columns have the following meaning (figure
\ref{objects_tree_columns}):

\begin{description}
	\item[Icon] a small picture identifying the node: the thumbnail of the formula for fractals,
	the shape icon for primitives (box, sphere, \ldots) and a colored symbol for groups.
	\item[Name] the display name of the node. Fractal nodes are automatically renamed to the name
	of the formula currently selected for them. To give a node your own name, right-click it and
 choose \emph{Rename} from the context menu (figure \ref{objects_tree_context_menu}).
	\item[Type] the kind of the node: \emph{fractal}, \emph{primitive}, \emph{hybrid},
	\emph{boolean Add}, \emph{boolean Mul} or \emph{boolean Sub}.
	\item[Material] a miniature preview of the material assigned to the object. Left-clicking the
	miniature opens the material manager, in which another material can be picked (chapter
	\ref{materials}). See section \ref{objects-tree-materials} below.
	\item[Enable] a checkbox which includes or excludes the object from rendering. Unchecking it
	``hides'' the object (and its whole subtree) without deleting it.
	\item[ID] the number of the node. All parameters of the node in the settings file are stored
	under names starting with \texttt{node\_XXXX\_} (for example \texttt{node\_0001\_position}),
	where \texttt{XXXX} is this number.
	\item[Object ID] the number of the underlying object (fractal slot, primitive or group). It is
	an internal identifier used mainly for debugging and when editing settings files by hand.
\end{description}

\simpleImageWithCaption75Width{img/manual/media/objects_tree_columns.png}
{The tree columns: Icon, Name, Type, Material, Enable, ID, Object ID}
{objects_tree_columns}{H}

\subsection{Adding objects}\label{objects-tree-adding}\index{objects tree!adding objects}

\subsubsection{Adding a fractal}

Pressing \emph{Add Fractal} creates a new fractal node (initially the Mandelbulb formula). The
formula itself and all its dedicated parameters (power, folds, etc.) are chosen in the editor panel
below the tree; the icon and the name of the node update automatically when the formula is changed.
Every fractal node has its own, completely independent set of formula parameters, so the same
formula can appear many times in one scene with different settings.

\subsubsection{Adding a primitive}

\emph{Add Primitive} opens a dialog in which the type of the shape can be chosen (figure
\ref{objects_tree_primitive_dialog}). Hovering over an icon shows a short description of the shape.
The available primitives are: rectangle, circle, box, cylinder, cone, sphere, torus, plane, water
(a wave-displaced surface for rendering seas and oceans), prism and ellipsoid. Primitives have
simple parameters (sizes, radii, angles) and are very fast to render, which makes them ideal for
``real world'' elements around a fractal: ground planes, walls, mirror boxes, etc.

\simpleImageWithCaptionHalfWidth{img/manual/media/objects_tree_primitive_dialog.png}
{The primitive type selector with a description tooltip}
{objects_tree_primitive_dialog}{H}

\subsubsection{Adding a group}

\emph{Add group} opens the ``Select group type'' dialog (figure \ref{objects_tree_group_dialog}) with
four kinds of container:

\begin{description}
	\item[Hybrid fractal] the formulas of all child fractals are combined into \emph{one} iteration
	sequence: the orbit is calculated by applying the formulas one after another according to their
	iteration counts. This is the classic Mandelbulber ``hybrid'' mode (see chapter
	\ref{iteration-loop} for the mathematics). A hybrid group accepts only fractal children.
	\item[Boolean ADD] (union) merges all children into a single object; the result shows every
	point contained in any of the children.
	\item[Boolean MUL] (intersection) keeps only the volume that lies \emph{inside all} children at
	the same time.
	\item[Boolean SUB] (difference) cuts the volume of the following children out of the first
	child (A $-$ B $-$ C \ldots). The order of the children matters: use drag and drop to change
	it (section \ref{objects-tree-drag-drop}).
\end{description}

\simpleImageWithCaptionThirdWidth{img/manual/media/objects_tree_group_dialog.png}
{Choosing the type of a new group}
{objects_tree_group_dialog}{H}

Groups are containers: while a group (or any of its children) is selected, the \emph{Add} buttons
put the new node \emph{inside} that group. Otherwise the new node is added at the scene root, under
\emph{World}.

A boolean group can be converted to another operator later without losing its children: right-click
the group and choose \emph{Change boolean operator}.

\subsection{Editing an object}\label{objects-tree-editor}\index{objects tree!editor panel}

Selecting a node in the tree builds the editor panel below it. The panel always starts with an
info line ``Editing: \ldots\ Type: \ldots'' showing the name and the type of the object that is
currently edited. The rest depends on the node type.

\subsubsection{Fractal nodes}

For a fractal the editor shows three blocks (figures
\ref{objects_tree_fractal_editor_formula}, \ref{objects_tree_general_params} and
\ref{objects_tree_calc_params}):

\begin{itemize}
	\item the \textbf{formula selector} and the parameters of the chosen formula (group boxes such
	as power, trigonometric functions, folding, etc.) -- the same parameters as in the older
	``fractal control'' interface, now per node; right next to the formula selector there is also
	the \emph{Don't add global C constant} checkbox, described below;
	\item the \textbf{General object parameters} (section \ref{objects-tree-general-params});
	\item the \textbf{Calculation parameters} (section \ref{objects-tree-calc-params}).
\end{itemize}

Because the formula part of the editor is usually taller than the visible pane, the panel is
presented in fragments: the figure below shows its top part, the two remaining blocks are shown
in figures \ref{objects_tree_general_params} and \ref{objects_tree_calc_params}.

\simpleImageWithCaptionHalfWidth{img/manual/media/objects_tree_fractal_editor_formula.png}
{The editor of a fractal node: formula selector and formula parameters}
{objects_tree_fractal_editor_formula}{H}

A fractal which is a child of a \emph{hybrid} group shows, in addition to the formula parameters,
a set of per-formula sequence controls (figure \ref{objects_tree_hybrid_children}) which define
the role that this formula plays in the shared iteration loop: \emph{Iterations}, \emph{Weight},
\emph{Start at iteration} / \emph{Stop at iteration}, \emph{Check for bailout condition} and
\emph{Don't add global C constant}. These options are not very intuitive, so each of them is
described in detail below.

\begin{description}
	\item[Iterations] In a hybrid group the orbit is not calculated by one formula, but is passed
	from formula to formula in blocks: the first formula runs this many consecutive iterations,
	then the next one runs its own block, and so on, cycling through the children's list until
	the global iteration budget (from the Calculation parameters of the group) is used up. This
	value is the length of the block dedicated to the formula. The default 1 makes the formulas
	alternate on every iteration; giving one formula clearly more iterations strengthens its
	shape in the final result at the cost of the others.
	\item[Weight] Tells how much of what the formula produced during one iteration is kept. At
	the end of every iteration the new orbit state is mixed with the state from before the
	iteration, using this value as the proportion. 1.0 (the default) means the result is fully
	taken over; 0.0 means the formula's changes are thrown away and the orbit passes through it
	unchanged; values in between blend the shapes, which is the main tool for morphing one
	fractal into another (see also the \emph{Formula weight (blending)} section in chapter
	\ref{iteration-loop}).
	\item[Start at iteration / Stop at iteration] The window of global iteration numbers in which
	this formula is active. Outside the window the formula is skipped completely, no matter what
	its block size is. A typical use is to let one formula work only at the beginning of the
	loop (to build the large-scale structure) and another one only towards the end (to add fine
	surface detail).
	\item[Check for bailout condition] The \emph{bailout} test is the condition which decides
	that a point does not belong to the fractal: after every iteration the length of the orbit
	vector is compared with a threshold (either the optimal value defined for the formula, or
	the common \emph{bailout} parameter from the Rendering engine pane). Once the orbit grows
	beyond the threshold the point has ``escaped'', the loop stops there and the point is
	colored as the outside area. This is what produces the classic fractal surface --- and it
	also saves time, because escaped pixels are not iterated further.
	With this option turned \emph{off}, no escape check is performed while this formula runs:
	the orbit is always iterated to the end of its window and such points are counted as
	\emph{inside} the set. In practice the formula is then rendered as a solid, filled shape,
	without interiors hollowed out by escaping orbits --- useful when the shape should act as a
	``filler'' or carrier detail inside the hybrid. The price is speed: pixels cannot bail out
	early during this part of the sequence, so every point consumes the full iteration budget.
	For a single, non-hybrid fractal the check is always performed, which is why the control is
	hidden outside hybrid groups.
	\item[Don't add global C constant] Classical Mandelbrot-type formulas add a constant to the
	iterated point after each transformation, for example $z_{n+1} = z_{n}^{2} + c$, where $c$ is
	a copy of the original point coordinates (or of the Julia constant in Julia mode). This
	addition is what makes the fractal grow from the point's position and keeps the shape
	attached to its source coordinates. Checking the box disables the addition for this formula
	only: the formula then iterates the orbit alone, in a Julia-like way. This checkbox is
	visible for every fractal node, and in a hybrid sequence it is usually desirable to disable
	it for the second and following formulas --- re-adding the original pixel point would keep
	dragging the orbit back towards its start and would distort the shape built by the first
	formula. The exact behavior depends on the formula: some add $c$ inside their own code, so
	the option is shown greyed out (not editable), and for some transform formulas the logic is
	inverted and the option appears as \emph{Add global C constant} (disabled addition by
	default). When the result suddenly shrinks to a small blob or disappears, toggling this
	option is the first thing to try.
\end{description}

Such a child does not carry its own material or transform -- the whole hybrid sequence shares
the parameters of the parent group node.

\simpleImageWithCaptionThirdWidth{img/manual/media/objects_tree_hybrid_children.png}
{A child fractal of a hybrid group with its per-formula mixing settings}
{objects_tree_hybrid_children}{H}

\subsubsection{Primitive nodes}

For a primitive the editor shows the parameters specific to its shape (for example the three side
lengths and corner rounding of a box) followed by the General object parameters block (figure
\ref{objects_tree_primitive_editor}).

\simpleImageWithCaptionThirdWidth{img/manual/media/objects_tree_primitive_editor.png}
{The editor panel of a box primitive}
{objects_tree_primitive_editor}{H}

\subsubsection{Group nodes}

For a hybrid or boolean group the editor shows the General object parameters and the Calculation
parameters which are then shared by all children of the group. Boolean groups have no
\emph{repeat} (it makes sense only for individual objects).

\subsubsection{General object parameters}\label{objects-tree-general-params}
\index{objects tree!position}\index{objects tree!repeat}\index{objects tree!scale}

This block (figure \ref{objects_tree_general_params}) is common to every object:

\begin{description}
	\item[Material] the material assigned to the object. Materials are created and edited in the
	Materials dock (chapter \ref{materials}).
	\item[Place object by mouse pointer] switches the mouse-click function (in the Navigation dock)
	to ``Place object'' mode. The next click on the rendered image moves this object to the clicked
	point; the object can also be dragged across the image while the mode is active
	(see section \ref{objects-tree-visual-tools}).
	\item[position (x, y, z)] the location of the object center, expressed in the coordinate system
	of the \emph{parent} node (see section \ref{objects-tree-hierarchy}).
	\item[X/Y/Z-axis rotation] rotation of the object around the three axes, in degrees.
	\item[scale] the size of the object. 1.0 keeps the original size.
	\item[Repeat (x, y, z)] repeats the object in an unlimited spatial grid. The values are the
	distances between neighboring copies along each axis; 0 on all axes disables the repetition.
	Repetition happens in the local space of the object, so a repeated object can be rotated
	freely without breaking the tiling.
	\item[detail level multiplier] per-object scaling of the global \emph{detail level} from the
	Rendering engine pane. Values above 1.0 render this object with more precision (slower),
	values below 1.0 with less.
	\item[Smooth DE combine] when enabled, the surface of this object blends smoothly with the
	surfaces of the other objects combined with it by a Boolean ADD (union). The
	\emph{Combine distance} parameter controls the width of the melted zone.
\end{description}

\simpleImageWithCaptionThirdWidth{img/manual/media/objects_tree_general_params.png}
{General object parameters shared by every node}
{objects_tree_general_params}{H}

\subsubsection{Calculation parameters}\label{objects-tree-calc-params}
\index{objects tree!julia mode}

Fractal and group nodes also have a \emph{Calculation parameters} box (figure
\ref{objects_tree_calc_params}) controlling how the formula is iterated:

\begin{description}
	\item[Maximum number of fractal iterations] the iteration budget used for this object (or for
	the whole hybrid group).
	\item[Julia mode] replaces the per-pixel iteration constant $c$ with a fixed Julia constant,
	turning the formula into its Julia version (section \ref{julia-mode}).
	\item[Julia constant (c)] the constant used when Julia mode is enabled.
	\item[Constant multiplier] multiplies the constant $c$ added on every iteration; 1.0 means no
	scaling. A cheap way to distort or ``melt'' a fractal.
	\item[Initial w-axis value] the starting value of the fourth (w) coordinate, used by 4D and
	hypercomplex formulas.
\end{description}

\simpleImageWithCaptionThirdWidth{img/manual/media/objects_tree_calc_params.png}
{Calculation parameters of a node}
{objects_tree_calc_params}{H}

\subsection{Organizing the tree}\label{objects-tree-organizing}\index{objects tree!drag and drop}

\subsubsection{Moving objects by drag and drop}\label{objects-tree-drag-drop}

Any node can be dragged with the left mouse button and dropped:

\begin{itemize}
	\item \emph{onto a group} (the group becomes its new parent -- the drop target is highlighted);
	\item \emph{above or below another node} (a horizontal line indicates the position -- the node
	becomes a sibling, which is how the order of children inside a boolean group is changed).
\end{itemize}

Some moves are rejected automatically: an object cannot be dropped into itself or into one of its
own descendants, primitives cannot become children of a hybrid group, and \emph{World} itself is
not a drop target. Figure \ref{objects_tree_drag_drop} shows the drop indicator.

\simpleImageWithCaptionHalfWidth{img/manual/media/objects_tree_drag_drop.png}
{Re-parenting and reordering nodes with drag and drop}
{objects_tree_drag_drop}{H}

\subsubsection{Duplicate}

\emph{Duplicate} makes a complete copy of the selected node \emph{together with its whole subtree}
(all children, grandchildren, \ldots). The copy receives new node and object ids, all its parameters
are copied and it is inserted next to the original in the same group. This is the quickest way to
create variations: duplicate a group, then change one parameter of the copy.

\subsubsection{Delete object}

\emph{Delete object} removes the selected node from the tree (for primitives all parameters of the
shape are deleted too). Children of the deleted node are \emph{not} lost: they are moved one level
up and become children of the deleted node's parent. This makes it possible to ``dissolve'' a group
without deleting the objects it contains. The \emph{World} root cannot be deleted.

\subsubsection{Context menu}

Right-clicking a node opens a small menu with \emph{Rename} and -- for boolean groups only --
\emph{Change boolean operator} (figure \ref{objects_tree_context_menu}).

\simpleImageWithCaptionThirdWidth{img/manual/media/objects_tree_context_menu.png}
{The context menu of a boolean group}
{objects_tree_context_menu}{H}

\subsection{Assigning materials directly in the tree}\label{objects-tree-materials}
\index{objects tree!material}

The \emph{Material} column contains a miniature of the assigned material. A left-click on the
miniature opens the material manager (figure \ref{objects_tree_material_manager}), where any
material can be selected for this node without leaving the Objects tab; choosing ``no material''
clears the miniature.

Objects inside a \emph{hybrid} group have no own material -- the whole hybrid sequence is colored
by the material of the parent group node. \emph{Boolean} group nodes also usually have no material:
the visible surface belongs to the child objects, so each child carries its own material, and the
color of the union/intersection follows the child that owns the surface point.

\simpleImageWithCaptionHalfWidth{img/manual/media/objects_tree_material_manager.png}
{Picking a material from the tree with one click}
{objects_tree_material_manager}{H}

\subsection{Visual placement tools}\label{objects-tree-visual-tools}
\index{objects tree!wire-frame preview}\index{objects tree!place object}

Two helpers make it easier to arrange objects in the scene:

\begin{description}
	\item[Display wire-frame previews] this checkbox turns on an overlay drawn on top of the render
	window: all primitives are shown as wireframes, every object gets a small origin cross (a 3D
	axes gizmo) and the currently selected node is highlighted. The overlay updates automatically
	while parameters are edited, so objects can be positioned ``by eye'' even before rendering
	(figure \ref{objects_tree_wireframe}).
	\item[Place object by mouse pointer] (the button inside the General object parameters block)
	switches the mouse-click function of the Navigation dock to ``Place object'' mode (the combo box
	entry lists the name of the selected node). The
	next click on the rendered image sets the position of this object to the clicked point (at the
	depth of the visible surface), and pressing and dragging the mouse button moves the object
	along the surface. Works for fractals, primitives and groups alike.
\end{description}

\simpleImageWithCaption75Width{img/manual/media/objects_tree_wireframe.png}
{Wire-frame previews and origin crosses help to position objects}
{objects_tree_wireframe}{H}

\subsection{The hierarchy and transforms}\label{objects-tree-hierarchy}
\index{objects tree!hierarchy}

Node transforms are \emph{inherited}: the position, rotation and scale of a child are applied
inside the coordinate system of its parent group. In practice this means:

\begin{itemize}
	\item moving or rotating a group moves \emph{all} of its children together, keeping their
	relative arrangement -- a convenient way to work on a sub-scene;
	\item scaling a group scales everything inside it;
	\item the \emph{repeat} of a leaf object is applied in the object's own local space, after the
	inherited transformation.
\end{itemize}

The world-space position shown in the status overlays and used by ``Place object'' is converted
automatically to the parent coordinate system, so nested groups can be positioned intuitively.

\subsection{Worked example: carving a sphere out of a Menger sponge}\label{objects-tree-example}
\index{objects tree!example}

This short recipe demonstrates the typical workflow:

\begin{enumerate}
	\item In the Objects tree press \emph{Add Fractal}. With the new node selected, choose the
	\emph{Menger Sponge} formula in the editor and press Render -- a sponge appears.
	\item Press \emph{Add group} and choose \emph{Boolean SUB}. A new group node appears at the
	scene root.
	\item Drag the Menger Sponge node onto the group -- it becomes the \emph{first} child of the
	group.
	\item Select the group and press \emph{Add Primitive}, choose \emph{Sphere}. The sphere becomes
	the second child; remember that Boolean SUB cuts the following children out of the first one.
	\item Enable \emph{Display wire-frame previews}, press \emph{Place object by mouse pointer} and
	drag the sphere across the image until it cuts a hole in the corner of the sponge.
	\item Duplicate the whole group, scale the copy down and repeat it, or add more shapes to the
	group. Experiment with the \emph{Smooth DE combine} option on the children of a Boolean ADD
	group to melt shapes together.
\end{enumerate}

\simpleImageWithCaptionHalfWidth{img/manual/media/objects_tree_example_result.png}
{Boolean SUB in action: a sphere cut out of a Menger sponge}
{objects_tree_example_result}{H}

\subsection{Notes and tips}\label{objects-tree-tips}

\begin{itemize}
	\item The \emph{Refresh} button rebuilds the tree from the current parameters. Use it when the
	scene was modified somewhere else (for example after loading settings or by a script) and the
	list looks outdated.
	\item Disabling nodes with the \emph{Enable} column instead of deleting them allows quick
	comparisons of scene variants.
	\item Nodes, not objects, carry the transforms: if a fractal is used by several groups it must
	be duplicated -- each copy has its own position, rotation, scale and material.
	\item Old scenes which were created before the Objects tree (single fractal, legacy hybrid or
	boolean arrangements) are migrated to the tree automatically on load; the corresponding legacy
	parameters remain stored for backward compatibility with other versions.
	\item The formula selector is grouped by fractal family, which makes it easy to discover
	formulas (section \ref{fractal-families}).
\end{itemize}
